
\documentclass[utf8, 14pt]{beamer}


%----------------------------------------------------------------------------------------
%	PACKAGES AND THEMES
%----------------------------------------------------------------------------------------

\mode<presentation> {% ou article

% The Beamer class comes with a number of default slide themes
% which change the colors and layouts of slides. Below this is a list
% of all the themes, uncomment each in turn to see what they look like.

%\usetheme{default}
%\usetheme{AnnArbor}
%\usetheme{Antibes}
%\usetheme{Bergen}
%\usetheme{Berkeley}
%\usetheme{Berlin}
%\usetheme{Boadilla}
%\usetheme{CambridgeUS}
%\usetheme{Copenhagen}
%\usetheme{Darmstadt}
%\usetheme{Dresden}
%\usetheme{Frankfurt}
%\usetheme{Goettingen}
%\usetheme{Hannover}
%\usetheme{Ilmenau}
%\usetheme{JuanLesPins}
%\usetheme{Luebeck}
\usetheme{Madrid}
%\usetheme{Malmoe}
%\usetheme{Marburg}
%\usetheme{Montpellier}
%\usetheme{PaloAlto}
%\usetheme{Pittsburgh}
%\usetheme{Rochester}
%\usetheme{Singapore}
%\usetheme{Szeged}
%\usetheme{Warsaw}

% As well as themes, the Beamer class has a number of color themes
% for any slide theme. Uncomment each of these in turn to see how it
% changes the colors of your current slide theme.

%\usecolortheme{albatross}
%\usecolortheme{beaver}
%\usecolortheme{beetle}
%\usecolortheme{crane}
%\usecolortheme{dolphin}
%\usecolortheme{dove}
%\usecolortheme{fly}
%\usecolortheme{lily}
%\usecolortheme{orchid}
%\usecolortheme{rose}
%\usecolortheme{seagull}
%\usecolortheme{seahorse}
%\usecolortheme{whale}
%\usecolortheme{wolverine}

%\setbeamertemplate{footline} % To remove the footer line in all slides uncomment this line

%\setbeamertemplate{footline}[page number] % To replace the footer line in all slides with a simple slide count uncomment this line

\setbeamertemplate{navigation symbols}{} % To remove the navigation symbols from the bottom of all slides uncomment this line
}


%\usepackage[utf8]{inputenc}
\usepackage[french]{babel}
\usepackage[T1]{fontenc}
\usepackage{lmodern}

%\usepackage{multimedia} %\movie[(options)]{image}{(lien vidéo)}
\usepackage{graphicx} % Allows including images
\usepackage{booktabs} % Allows the use of \toprule, \midrule and \bottomrule in tables
\usepackage{eurosym}
\usepackage{multicol}
\usepackage{tikz-uml}

\graphicspath{{../images/}}

% ------------------------------------------------------ %

%\defbeamertemplate*{footline}{shadow theme}
%{%
%  \leavevmode%
%  \hbox{\begin{beamercolorbox}[wd=.5\paperwidth,ht=2.5ex,dp=1.125ex,leftskip=.3cm plus1fil,rightskip=.3cm]{author in head/foot}%
%    \usebeamerfont{author in head/foot}\insertframenumber\,/\,\inserttotalframenumber\hfill\insertshortauthor
%  \end{beamercolorbox}%
%  \begin{beamercolorbox}[wd=.5\paperwidth,ht=2.5ex,dp=1.125ex,leftskip=.3cm,rightskip=.3cm plus1fil]{title in head/foot}%
%    \usebeamerfont{title in head/foot}\insertshorttitle%
%  \end{beamercolorbox}}%
%  \vskip0pt%
%}

% ------------------------------------------------------ %


%----------------------------------------------------------------------------------------
%	TITLE PAGE
%----------------------------------------------------------------------------------------

\title[Design Patterns]{Design Patterns}
\subtitle{OS Virtuel}

\author[Axel G. | Audoin De C.]{Axel \textsc{Gendillard} \and Audoin \textsc{De Chantérac}}
\institute[]
{
	Étudiant ingénieur\\ 
	Groupe ESEO \\
	%\medskip
	%\textit{axel.gendillard@reseau.eseo.fr}
}

\date{\today} 

%\logo{\includegraphics[scale=0.3]{logo}}

\titlegraphic{%
%\includegraphics[scale=0.55]{logo} \hspace{1cm}~%
%\includegraphics[scale=0.6]{eseo}
}


\begin{document}

\begin{frame}
\titlepage
\end{frame}

\begin{frame}{Sommaire}
\tableofcontents
\end{frame} 

%----------------------------------------------------------------------------------------
%	PRESENTATION SLIDES
%----------------------------------------------------------------------------------------

\section{Incrément 1 : Singleton}
	
	\begin{frame}{Incrément 1 : Singleton}
		\begin{block}{Objectifs}
		\begin{itemize}
		\item Créer un nouvel utilisateur à chaque validation;
		\item Avoir un utilisateur unique par login;
		\item Créer plusieurs terminaux par utilisateur;
		\item Interdir la création d'utilisateur hors de la fenêtre.
		\end{itemize}
		\end{block}
		
		\begin{exampleblock}{Patron Singleton}
		\begin{itemize}
		\item Garantir l'unicité de l'instance d'une classe;
		\item Accéder globalement à l'instance.
		\end{itemize}
		\end{exampleblock}					
	\end{frame}

	\begin{frame}{Incrément 1 : Singleton - UML}
		\begin{figure}[!h]
		\centering
		\includegraphics[width=\textwidth]{../uml/uml-singleton}
		\end{figure}		
	\end{frame}


\section{Incrément 2 : Observer}
	
	\begin{frame}{Incrément 2 : Observer}
		\begin{block}{Objectifs}
		\begin{itemize}
		\item Synchroniser les historiques des terminaux;
		\item Mettre à jour tous les terminaux lors de l'entrée d'une commande.
		\end{itemize}
		\end{block}
		
		\begin{exampleblock}{Patron Observer}
		\begin{itemize}
		\item Permet de mettre à jour les objets observés en fonction de l'objet observateur.
		\end{itemize}
		\end{exampleblock}					
	\end{frame}

	\begin{frame}{Incrément 2 : Observer - UML}
		\begin{figure}[!h]
		\centering
		\includegraphics[width=\textwidth]{../uml/uml-observer}
		\end{figure}		
	\end{frame}



\section{Incrément 3 : Visiteur}
	
	\begin{frame}{Incrément 3 : Visiteur}
		\begin{block}{Objectifs}
		\begin{itemize}
		\item Interpréter des commandes du terminal (ls et cat) pour chaque type de noeud;
		\item Utiliser le patron Visitor.
		\end{itemize}
		\end{block}
		
		\begin{exampleblock}{Patron Visiteur}
		\begin{itemize}
		\item Sépare l'algorithme de la structure de données;
		\item Chaque élément implémente une méthode d'acceptation de visiteur;
		\item Chaque visiteur implémente une méthode \emph{visit()} pour chaque élément (file, folder et link).
		\end{itemize}
		\end{exampleblock}					
	\end{frame}

	\begin{frame}{Incrément 3 : Visiteur - UML}
		\begin{figure}[!h]
		\centering
		\includegraphics[width=\textwidth]{../uml/uml-visitor}
		\end{figure}		
	\end{frame}


\section{Incrément 4 : Commande}
	
	\begin{frame}{Incrément 4 : Commande}
		\begin{block}{Objectifs}
		\begin{itemize}
		\item Réutiliser les traitements interprétés par le terminal;
		\item Exprimer de manière unique le traitement des commandes ls et cat.
		\end{itemize}
		\end{block}
		
		\begin{exampleblock}{Patron Commande}
		\begin{itemize}
		\item Isole une requête;
		\item Ces requêtes peuvent provenir de plusieurs émetteurs (\emph{TerminalOS} et \emph{ExplorerOS});
		\item Ces émetteurs doivent produire la même requête;
		\item Ces requêtes doivent pouvoir être annulées.
		\end{itemize}
		\end{exampleblock}					
	\end{frame}

	\begin{frame}{Incrément 4 : Commande - UML}
		\begin{figure}[!h]
		\centering
		\includegraphics[width=0.78\textwidth]{../uml/uml-command}
		\end{figure}		
	\end{frame}

\section{Incrément 5 : Commande ++}
	
	\begin{frame}{Incrément 5 : Commande ++}
		\begin{block}{Objectifs}
		\begin{itemize}
		\item RM : Supprimer un dossier, un fichier ou un lien;
		\item MKDIR : Créer un dossier;
		\item TOUCH : Créer un fichier;
		\item LN : Créer un lien.
		\end{itemize}
		\end{block}
		
		\begin{exampleblock}{Patron Commande}
		\begin{itemize}
		\item Isole une requête;
		\item Ces requêtes peuvent provenir de plusieurs émetteurs (\emph{TerminalOS} et \emph{ExplorerOS}) ayant le même focntionnement;
		\item Ces requêtes doivent pouvoir être annulées.
		\end{itemize}
		\end{exampleblock}					
	\end{frame}

	\begin{frame}{Incrément 5 : Commande - UML}
		\begin{figure}[!h]
		\centering
		\includegraphics[width=0.78\textwidth]{../uml/uml-command}
		\end{figure}		
	\end{frame}

\section{Incrément 6 : Proxy}
	
	\begin{frame}{Incrément 6 : Proxy}
		\begin{block}{Objectifs}
		\begin{itemize}
		\item Lecture seule : Pas de modification des données des utilisateurs ls et cat sont autorisées;
		\item Lecture/écriture : Les commandes mkdir, touch, ln et rm sont autorisées.
		\end{itemize}
		\end{block}
		
		\begin{exampleblock}{Patron Proxy}
		\begin{itemize}
		\item Rôle d'aiguilleur en fonction des droits de l'utilisateur;
		\item Couche d'abstraction entre le client et la commande.
		\end{itemize}
		\end{exampleblock}					
	\end{frame}

	\begin{frame}{Incrément 6 : Proxy - UML}
		\begin{figure}[!h]
		\centering
		\includegraphics[width=\textwidth]{../uml/uml-proxy}
		\end{figure}		
	\end{frame}


\end{document} 
